\documentclass[10pt,a4paper]{amsart}
\usepackage{amsmath,amssymb,amsthm,graphicx}
% Retained minimal preamble for the Stacks Project chapter more-algebra.tex (More on Algebra).
% The chapter file pinned-original-source.tex carries only \input{preamble} on its first line, so no
% preamble text lives inside the chapter itself.  The Stacks Project uses one single shared preamble for
% every chapter; the snapshot stages it identically as sources/stacks-chow/preamble.tex and
% sources/stacks-etale-cohomology/preamble.tex.  The block below is copied mechanically, line for line,
% from sources/stacks-chow/preamble.tex lines 14-18 (the Xy-pic setup) and 154-217 (the native theorem
% declarations together with the mathematical macros).  Nothing was added, removed or redefined; the
% wrapper class and the amsmath/amssymb/amsthm/graphicx base packages are supplied by the assembler.
% Chapter: more-algebra.tex -- More on Algebra.
% For commutative diagrams we use Xy-pic
\usepackage[all]{xy}

% We use 2cell for 2-commutative diagrams.
\xyoption{2cell}

% Theorem environments.
%
\theoremstyle{plain}
\newtheorem{theorem}[subsection]{Theorem}
\newtheorem{proposition}[subsection]{Proposition}
\newtheorem{lemma}[subsection]{Lemma}

\theoremstyle{definition}
\newtheorem{definition}[subsection]{Definition}
\newtheorem{example}[subsection]{Example}
\newtheorem{exercise}[subsection]{Exercise}
\newtheorem{situation}[subsection]{Situation}

\theoremstyle{remark}
\newtheorem{remark}[subsection]{Remark}
\newtheorem{remarks}[subsection]{Remarks}

\numberwithin{equation}{subsection}

% Macros
%
\def\lim{\mathop{\mathrm{lim}}\nolimits}
\def\colim{\mathop{\mathrm{colim}}\nolimits}
\def\Spec{\mathop{\mathrm{Spec}}}
\def\Hom{\mathop{\mathrm{Hom}}\nolimits}
\def\Ext{\mathop{\mathrm{Ext}}\nolimits}
\def\SheafHom{\mathop{\mathcal{H}\!\mathit{om}}\nolimits}
\def\SheafExt{\mathop{\mathcal{E}\!\mathit{xt}}\nolimits}
\def\Sch{\mathit{Sch}}
\def\Mor{\mathop{\mathrm{Mor}}\nolimits}
\def\Ob{\mathop{\mathrm{Ob}}\nolimits}
\def\Sh{\mathop{\mathit{Sh}}\nolimits}
\def\NL{\mathop{N\!L}\nolimits}
\def\CH{\mathop{\mathrm{CH}}\nolimits}
\def\proetale{{pro\text{-}\acute{e}tale}}
\def\etale{{\acute{e}tale}}
\def\QCoh{\mathit{QCoh}}
\def\Ker{\mathop{\mathrm{Ker}}}
\def\Im{\mathop{\mathrm{Im}}}
\def\Coker{\mathop{\mathrm{Coker}}}
\def\Coim{\mathop{\mathrm{Coim}}}

% Boxtimes
%
\DeclareMathSymbol{\boxtimes}{\mathbin}{AMSa}{"02}

%
% Macros for moduli stacks/spaces
%
\def\QCohstack{\mathcal{QC}\!\mathit{oh}}
\def\Cohstack{\mathcal{C}\!\mathit{oh}}
\def\Spacesstack{\mathcal{S}\!\mathit{paces}}
\def\Quotfunctor{\mathrm{Quot}}
\def\Hilbfunctor{\mathrm{Hilb}}
\def\Curvesstack{\mathcal{C}\!\mathit{urves}}
\def\Polarizedstack{\mathcal{P}\!\mathit{olarized}}
\def\Complexesstack{\mathcal{C}\!\mathit{omplexes}}
% \Pic is the operator that assigns to X its picard group, usage \Pic(X)
% \Picardstack_{X/B} denotes the Picard stack of X over B
% \Picardfunctor_{X/B} denotes the Picard functor of X over B
\def\Pic{\mathop{\mathrm{Pic}}\nolimits}
\def\Picardstack{\mathcal{P}\!\mathit{ic}}
\def\Picardfunctor{\mathrm{Pic}}
\def\Deformationcategory{\mathcal{D}\!\mathit{ef}}

% --- appended 2026-10-12 (Chief): the project own definitions of the reference/slogan/history
% display environments, copied mechanically from the shared Stacks preamble lines 8-12
% (\usepackage{verbatim} plus the three \newenvironment lines, which the chapter preamble keeps
% verbatim from the project preamble). Nothing added, removed or redefined: these environments
% carry no mathematical content and the project itself typesets them as comments. Without these
% three lines a counted statement whose own text uses them fails with Environment slogan undefined.
\usepackage{verbatim}
\newenvironment{reference}{\comment}{\endcomment}
\newenvironment{slogan}{\comment}{\endcomment}
\newenvironment{history}{\comment}{\endcomment}
\usepackage{cleveref}
\crefname{originallemma}{original lemma}{original lemmas}
\Crefname{originallemma}{Original lemma}{Original lemmas}
\makeatletter
\newlabel{lemma-get-K-flat-system}{{091G}{source}{}{sourcepointer-1}{}}
\newlabel{lemma-get-K-flat-system@cref}{{[originallemma][0][]091G}{[1][0][]source}{}{}{}}
\makeatother
\crefname{originallemma}{original lemma}{original lemmas}
\Crefname{originallemma}{Original lemma}{Original lemmas}
\makeatletter
\newlabel{lemma-K-flat-quasi-isomorphism}{{06Y0}{source}{}{sourcepointer-2}{}}
\newlabel{lemma-K-flat-quasi-isomorphism@cref}{{[originallemma][0][]06Y0}{[1][0][]source}{}{}{}}
\makeatother
\crefname{originallemma}{original lemma}{original lemmas}
\Crefname{originallemma}{Original lemma}{Original lemmas}
\makeatletter
\newlabel{lemma-derived-tor-quasi-isomorphism-other-side}{{064L}{source}{}{sourcepointer-3}{}}
\newlabel{lemma-derived-tor-quasi-isomorphism-other-side@cref}{{[originallemma][0][]064L}{[1][0][]source}{}{}{}}
\makeatother
\crefname{originallemma}{original lemma}{original lemmas}
\Crefname{originallemma}{Original lemma}{Original lemmas}
\makeatletter
\newlabel{lemma-derived-tor-exact}{{064J}{source}{}{sourcepointer-4}{}}
\newlabel{lemma-derived-tor-exact@cref}{{[originallemma][0][]064J}{[1][0][]source}{}{}{}}
\makeatother
\title{Derived tensor products for systems of rings and their functoriality}
\author{The Stacks Project Authors}
\date{stacks-project:more-algebra.tex:local-snapshot-20261012; GNU FDL 1.2 or later}
\begin{document}
\maketitle
\setcounter{subsection}{1}
\noindent\textit{Editorial note.} All delivered mathematical text is the author's own, copied from the pinned Stacks Project chapter file more-algebra.tex; no statement or proof step was rewritten, reordered or removed. The counted span is the author environment and its own proof environment. References inside the span are handed to the assembler as navigation stubs that carry the canonical tag of the cited tag only; whether each of them is load-bearing is recorded on the stub itself.\par\medskip
\noindent\textit{Allowed standard prerequisites (editorial).}\begin{itemize}
\item Every object of the derived category of systems admits a K-flat representative system. invoked at proof lines 23869-23870. Provides the representative for K.
\item A K-flat representative can be used to compute derived tensor products, and the result is independent of the representative on the other side, by the K-flat quasi-isomorphism and derived-tor lemmas. invoked at proof lines 23877-23879. Well-definedness in both variables.
\item Exactness of the termwise derived tensor in each variable. invoked at proof lines 23881-23882. Gives the triangulated functoriality.
\end{itemize}\medskip
\ifdefined\hypertarget\hypertarget{sourcepointer-1}{}\fi
\noindent\textit{Source reference.} Original lemma 091G: sources/stacks-more-algebra/pinned-original-source.tex, source lines 23816--23816. Cited at proof lines 23870-23870 for the representative of K; the definition and well-definedness are inside the counted span.\par\medskip
\ifdefined\hypertarget\hypertarget{sourcepointer-2}{}\fi
\noindent\textit{Source reference.} Original lemma 06Y0: sources/stacks-more-algebra/pinned-original-source.tex, source lines 14192--14192. Cited at proof lines 23878-23878 for well-definedness; the conclusion is inside the counted span.\par\medskip
\ifdefined\hypertarget\hypertarget{sourcepointer-3}{}\fi
\noindent\textit{Source reference.} Original lemma 064L: sources/stacks-more-algebra/pinned-original-source.tex, source lines 14457--14457. Cited at proof lines 23878-23879 for well-definedness in the second variable; the conclusion is inside the counted span.\par\medskip
\ifdefined\hypertarget\hypertarget{sourcepointer-4}{}\fi
\noindent\textit{Source reference.} Original lemma 064J: sources/stacks-more-algebra/pinned-original-source.tex, source lines 14153--14153. Cited at proof lines 23881-23882 as the model for the exactness argument; the definition and well-definedness are inside the counted span.\par\medskip
\begin{lemma}
\label{lemma-derived-tensor-product-systems}
Let $(A_n)$ be an inverse system of rings. Given
$K, L \in D(\textit{Mod}(\mathbf{N}, (A_n)))$ there is a canonical derived
tensor product $K \otimes^\mathbf{L} L$ in $D(\mathbf{N}, (A_n))$
compatible with the maps to $D(A_n)$. The construction is symmetric
in $K$ and $L$ and an exact functor of triangulated categories in
each variable.
\end{lemma}

\begin{proof}
Choose a representative $(K_n^\bullet)$ for $K$ such that each $K_n^\bullet$
is a K-flat complex (Lemma \ref{lemma-get-K-flat-system}).
Then you can define $K \otimes^\mathbf{L} L$ as the object represented by
the system of complexes
$$
(\text{Tot}(K_n^\bullet \otimes_{A_n} L_n^\bullet))
$$
for any choice of representative $(L_n^\bullet)$ for $L$.
This is well defined in both variables by
Lemmas \ref{lemma-K-flat-quasi-isomorphism} and
\ref{lemma-derived-tor-quasi-isomorphism-other-side}.
Compatibility with the map to $D(A_n)$ is clear.
Exactness follows exactly as in
Lemma \ref{lemma-derived-tor-exact}.
\end{proof}

\end{document}
